\begin{resumo}[Abstract]
 \begin{otherlanguage*}{english}
    Education is a fundamental pillar for a country's development, but the Brazilian educational landscape faces challenges such as low motivation and student disengagement, reflected in indicators like Ideb and PISA. Collaborative learning methodologies, such as tutoring, and the application of gamification have shown great potential to mitigate these problems, increasing engagement and the effectiveness of the teaching-learning process. 
    The objective of this study was to develop a web application that mediates contact between tutors in specific fields of knowledge and individuals seeking tutoring. The platform provides features to search for and contact tutors, request sessions, and evaluate both tutors and learners after each interaction, utilizing gamification elements to foster engagement and consistent use of the application. It is an applied research, which uses bibliographic research procedures and agile software development methodologies. In the end of this work, a web application was developed and is currently available for use through standard web browsers.
   \vspace{\onelineskip}
 
   \noindent 
   \textbf{Key-words}: Tutoring; Recommendation System; Gamification; Web Application.
 \end{otherlanguage*}
\end{resumo}
